\newpage
\section{Auxiliary Functions for Synthesis}\label{sec:tech:algo}

This section describes three auxiliary functions used for generator
synthesis. The first of these, $\normPlan$, converts a SRE formula into a set of unsafe abstract traces. The second, $\mathbf{NormCtx}$ normalizes the \PAT{} of each effect operator. The third, $\termDerive$, generates a generator program from a refined abstract trace.

\begin{algorithm}[t!]
  \Procedure{$\normPlan(A) := $}{
  \Match{$A$}{
    \lCase{$\emptyset$}{ \Return{$\emptyset$} }
    \lCase{$\epsilon$}{ \Return{$\{\epsilon\}$} }
    \lCase{$\msg{op}{\phi}$}{ \Return{$\{\msg{op}{\phi}\}$} }
    \lCase{$A_1 \seqA A_2$}{ \Return{$\{ \pi_1 \seqA \Pi_2 ~|~ \Pi_1 \in \normPlan(A_1) \land \Pi_2 \in \normPlan(A_2) \}$}}
    \lCase{$A_1 \lorA A_2$}{ \Return{$\normPlan(A_1) \cup \normPlan(A_2)$}}
    \lCase{$A^*$}{\Return{$\{ A^* \}$}}
  }
  }
  \Procedure{$\mathbf{NormCtx}(\Delta) := $}{
    $\Delta \leftarrow \emptyset$\;
    \ForEach{$(\Theta, \overline{z{:}b}\garr\overline{y{:}t}\sarr\rg{H}{\eff{op}(\overline{y})}{F}) \in \Delta(\eff{op})$}{
      \ForEach{$\Pi_h \in \normPlan(H)$ and $\Pi_f \in \normPlan(F)$}{
        $\Delta'(\eff{op}) \leftarrow \Delta'(\eff{op}) \cup \{(\Theta, \overline{z{:}b}\garr\overline{y{:}t}\sarr\rg{\Pi_h}{\eff{op}(\overline{y})}{\Pi_f})\}$\;
      }
    }
    \Return{$\Delta'$}\;
  }
  \caption{Abstract Trace Normalization}
  \label{algo:norm}
\end{algorithm}

\paragraph{Normalization} 
As shown in \autoref{algo:norm}, the function $\normPlan$ recursively translates the input
SRE into a set of abstract traces. It keeps the Kleene star operator $A^*$ unchanged and divides $A \lorA A'$ into multiple abstract traces. The function $\mathbf{NormCtx}$ also normalizes the history and prophecy regexes in \PAT{}.


% \begin{lemma}\label{lemma:traceAnd}[Abstract traces are closed under
%     conjunction] The conjunction ($\land$) of two abstract traces is also
%     an abstract trace.
% \end{lemma}

\begin{lemma}\label{lemma:norm-sound}[Normalization is sound ]
  The normalized result has the same denotation as the input SRE,
  that is, for all SRE $A$ and set of traces $\{\Pi_i\}$, {\small
    $\denot{A} = \bigcup_{i} \denot{\Pi_i}$}
\end{lemma}

\begin{algorithm}[t!]
    \Procedure{$\S{TermDerive}(\Gamma, \Pi) := $}{
      \Match{$\Gamma$}{
      \Case{$[]$}{ \Return{$\S{DeriveTrace}(\Pi)$}\; }
      \Case{$x{:}\rawnuot{b}{\phi} \cons \Gamma'$}{
      \Return{$\zassume{\overline{x}}{\phi[\vnu \mapsto x]}{\S{TermDerive}(\Gamma', \Code{ToList}(\Pi))}$}\;
      }
    }
    }
    \caption{Term Derivation}
    \label{algo:term-derive}
  \end{algorithm}

\begin{algorithm}[t!]
    \Procedure{$\S{DeriveTrace}(\Pi) := $}{
    \Match{$\Pi$}{
      \lCase{$[]$}{ \Return{$()$} }
      \Case{$\Pi^* \cons \Pi'$}{ 
        $e \leftarrow \S{DeriveTrace}(\Pi)$\;
        $e' \leftarrow\S{DeriveTrace}(\Pi')$\;
        \Return{$\zwhile{e}{e'}$}
        }
      \Case{$\msgB{op}{\overline{x}}{\phi} \seqA \Pi'$ when $\mykw{gen}\;\eff{op}$}{
        $\overline{x'} \leftarrow \Code{GetFreshNames}(\overline{x})$\;
        $e \leftarrow \S{DeriveTrace}(\Pi')$\;
        \Return{$\zassume{\overline{x'}}{\phi\msubst{x}{x'}}{e}$}
      }
      \Case{$\msgB{op}{\overline{x}}{\phi} \cons \Pi'$ when $\mykw{obs}\;\eff{op}$}{
        $\overline{x'} \leftarrow \Code{GetFreshNames}(\overline{x})$\;
        $e \leftarrow \S{DeriveTrace}(\Pi')$\;
        \Return{$\zobs{\overline{x'}}{op}{}{\zassert{\phi\msubst{x}{x'}}{e}}$}
      }
    }
    }
    \caption{Trace Derivation}
    \label{algo:trace-derive}
\end{algorithm}

\paragraph{Term Derivation} The term derivation function
\(\S{TermDerive}\) is shown in \autoref{algo:term-derive}. It first
converts the input type context into \(\mykw{assume}\) statements over
the corresponding qualifiers in pure refinement types (line 6), then
derives the abstract trace with the help of the \(\S{DeriveTrace}\)
subroutine shown in \autoref{algo:trace-derive}. The input abstract
trace is first be converted into a list of SRE
(\(\Code{ToList}\)) before the subroutine is called; it then
recursively transforms this list into a generator program. The kleene star $\Pi^*$ is converted into a $\mykw{while}$ loop on line 4. For a generable symbolic event (line 8),
\(\S{DeriveTrace}\) inserts an \(\mykw{assume}\) expression before the
\(\mykw{gen}\) expression on line 11. Conversely, for observable
events, \(\S{DeriveTrace}\) adds an \(\mykw{assert}\) expression after
the \(\mykw{obs}\) expression on line 13.

\begin{lemma}\label{lemma:derive-sound}[Term Derivation is Sound]
  Fora given type context $\Gamma$, well-founded type context
  $\Delta$, abstract trace $\Pi$, and term $e$,
  {\small
    \begin{align*}
      \S{TermDerive}(\Gamma, \Pi) = e \impl \Gamma;\Delta;\emptyset \vdash e : \capTau{\emptyset, \rg{\epsilon}{A}{\epsilon}}
    \end{align*}}
\end{lemma}
